\chapter{Theoretical Framework}
\label{ch:theory}
% =========================================================================
\section{Introduction}
\label{sec:ch3-intro}

The empirical analysis of Chapters 6 and 7 documents a robust negative carbon-conditional cancellation differential between Blue and Green hydrogen projects: at low EUA prices, the cancellation hazard of Blue projects is substantially elevated relative to Green; at high EUA prices, the gap narrows or reverses. The interaction coefficient $\bint$ enters the hazard model in the form $\eta_{it} \supset \bint \cdot \mathrm{Blue}_i \cdot z_t$ and is estimated at approximately $-1.5$ in our cleanest data subsample.

This chapter develops a theoretical framework that derives such a coefficient from explicit project economics. The framework is consciously simple: we model a representative project sponsor making a cancellation decision in a real-options setting, with the EUA price as the single stochastic state variable. The simplicity is deliberate: a richer model with multiple state variables (gas prices, equipment costs, capital availability) would generate testable implications that are difficult to disentangle empirically. The single-state model isolates the carbon-pricing mechanism most cleanly, providing a benchmark against which empirical estimates can be compared.

Three contributions follow. First, we provide an explicit economic mechanism connecting the EUA price to the cancellation hazard differential: Blue projects bear higher initial capital sunk costs and retain a small residual operational exposure to carbon prices, while Green projects have lower sunk costs and essentially no operational carbon exposure. Second, we derive the comparative statics that map the model parameters into the sign and magnitude of the empirical interaction coefficient. Third, we situate our empirical findings within the broader literature on irreversible investment under uncertainty, connecting to \cite{pindyck1991irreversibility, dixit1994investment, bolton2021carbon, odenweller2022probabilistic}.

The remainder of the chapter is organised as follows. Section \ref{sec:ch3-setup} establishes the model setup. Section \ref{sec:ch3-solution} solves the cancellation problem. Section \ref{sec:ch3-comparative} derives comparative statics and the connection to $\bint$. Section \ref{sec:ch3-extensions} discusses extensions and limitations. Section \ref{sec:ch3-conclusion} concludes.

% =========================================================================
\section{Model Setup}
\label{sec:ch3-setup}

\subsection{State variable: EUA price}

The EUA spot price $z_t$ evolves as a geometric Brownian motion under the risk-neutral measure:
\begin{equation}
\frac{dz_t}{z_t} = \mu \, dt + \sigma \, dW_t,
\end{equation}
where $\mu$ captures the risk-neutral drift (interpretable as the expected rate of carbon-price appreciation given the announced Phase IV cap trajectory and Market Stability Reserve dynamics) and $\sigma$ is the diffusion parameter (interpretable as the standard deviation of EUA log-returns). Empirically, $\sigma$ is large: realised volatility of EUA spot returns over 2018--2024 is approximately 50\% annualised, comparable to that of small-cap equity indices.

We work with the standardised state $z_t$ which can take both positive and negative values relative to the historical mean.

\subsection{Project profits}

Consider a representative project of technology $\tau \in \{B, G\}$ (Blue versus Green). The project is characterised by an initial capital cost $K^{\tau}$ (sunk if cancellation is delayed past the construction milestone), a fixed operating cost $c^{\tau}_{\mathrm{fix}}$, and a per-period profit flow that depends on the EUA price.

We posit a linear-in-$z$ structure for per-period profits:
\begin{align}
\PB(z_t) &= a^B + b^B z_t - c^{B}_{\mathrm{fix}}, \\
\PG(z_t) &= a^G + b^G z_t - c^{G}_{\mathrm{fix}},
\end{align}
where:
\begin{itemize}
    \item $a^{\tau}$ is the deterministic component (reflecting hydrogen prices, subsidy income, fixed contracts);
    \item $b^{\tau}$ is the sensitivity of profit to the EUA price.
\end{itemize}

The key economic assumption is:
\begin{equation}
b^B > 0 \quad \text{and} \quad b^G > 0 \quad \text{with} \quad b^B > b^G.
\label{eq:ch3-key-assumption}
\end{equation}

The interpretation: both technologies benefit from a high EUA price because high EUA prices either (a) raise the equilibrium hydrogen market price as carbon-intensive (grey) hydrogen exits the market, or (b) trigger CBAM/border-adjustment mechanisms that protect domestic clean hydrogen producers. However, Blue's benefit is larger ($b^B > b^G$) because Blue is the marginal carbon-friendly producer that displaces grey hydrogen: at high EUA, grey exits and Blue's market share expands disproportionately, while Green's benefit is mediated only through the equilibrium price effect.

This assumption deliberately reverses the naïve intuition that ``Blue has carbon costs, Green does not, so Green benefits more from high EUA.'' That reasoning treats Blue's residual carbon exposure as the dominant channel; we argue the dominant channel is competitive displacement of grey hydrogen, in which Blue is the more direct beneficiary. We discuss alternative parameterisations in Section \ref{sec:ch3-extensions}.

\subsection{Capital structure}

We assume Blue projects have higher capital intensity:
\begin{equation}
K^B > K^G.
\label{eq:ch3-capital}
\end{equation}

This is consistent with industry data: a typical SMR-with-CCS facility requires CAPEX of order \$ 1000--1500 per kg/day of H$_2$ production capacity, while a PEM electrolyser plant requires \$ 600--1100 per kg/day. The CAPEX gap reflects the additional cost of carbon capture, transport, and sequestration infrastructure.

For the cancellation decision, the relevant quantity is the \textit{unrecoverable} portion of capital --- the difference $K^B - K^G$ that has been sunk before the cancellation decision is made. We assume both technologies have a recoverable salvage value of $S^{\tau} \leq K^{\tau}$, and write $\hat{K}^{\tau} = K^{\tau} - S^{\tau}$ as the irreversible component.

\subsection{Cancellation as an option}

At each time $t$, the sponsor chooses whether to continue the project (and continue paying fixed operating costs) or cancel and recover the salvage value $S^{\tau}$. The continuation value at time $t$ is:
\begin{equation}
V^{\tau}(z_t) = \E_t \left[ \int_t^{\infty} e^{-r(s-t)} \Pi^{\tau}(z_s) \, ds \, \Big| \, z_t \right] - \hat{K}^{\tau},
\end{equation}
where $r$ is the discount rate, $\Pi^{\tau}$ is the per-period profit function, and $\hat{K}^{\tau}$ captures the remaining capital to be sunk. The cancellation decision is:
\begin{equation}
\text{Cancel project } \tau \text{ at time } t \quad \iff \quad V^{\tau}(z_t) < 0.
\end{equation}

Equivalently, project $\tau$ is cancelled when the EUA price falls below a critical threshold $\bar{z}^{\tau}$ at which the option to wait is optimally exercised.

% =========================================================================
\section{Solution: Optimal Cancellation Policy}
\label{sec:ch3-solution}

\subsection{Threshold characterisation}

The optimal cancellation rule is a threshold: project $\tau$ is cancelled whenever $z_t < \bar{z}^{\tau}$. Standard arguments \cite{dixit1994investment} yield:
\begin{equation}
\bar{z}^{\tau} = \frac{r}{b^{\tau}} \left[ c^{\tau}_{\mathrm{fix}} - a^{\tau} + \hat{K}^{\tau} (r - \mu) \right] \cdot \frac{\eta_2}{\eta_2 - 1},
\label{eq:ch3-threshold}
\end{equation}
where $\eta_2 > 1$ is the positive root of the fundamental quadratic $\frac{1}{2}\sigma^2 \eta(\eta-1) + \mu \eta - r = 0$. The factor $\eta_2 / (\eta_2 - 1) > 1$ is the standard ``option premium'' that exceeds the Marshallian (zero-NPV) threshold, reflecting the value of waiting under uncertainty.

\subsection{Comparing Blue and Green thresholds}

From equation \eqref{eq:ch3-threshold}, the ratio of cancellation thresholds is:
\begin{equation}
\frac{\bar{z}^B}{\bar{z}^G} = \frac{b^G}{b^B} \cdot \frac{c^{B}_{\mathrm{fix}} - a^B + \hat{K}^B(r - \mu)}{c^{G}_{\mathrm{fix}} - a^G + \hat{K}^G(r - \mu)}.
\end{equation}

Two effects work in opposite directions:
\begin{enumerate}
    \item Blue's higher EUA sensitivity ($b^B > b^G$) implies that Blue requires a \textit{lower} EUA threshold to remain viable (the factor $b^G/b^B < 1$). High EUA prices keep Blue profitable; Blue therefore cancels at lower $z$.
    \item Blue's higher sunk capital ($\hat{K}^B > \hat{K}^G$) raises the breakeven required to recover the remaining investment.
\end{enumerate}
The net direction of $\bar{z}^B$ versus $\bar{z}^G$ depends on the relative magnitudes. Under the parameter regime where the capital differential dominates (high $\hat{K}^B - \hat{K}^G$, modest $b^B/b^G$ differential), Blue cancels at a \textit{higher} EUA threshold than Green: $\bar{z}^B > \bar{z}^G$.

\subsection{Hazard rate implications}

Under the assumption that $z_t$ follows a GBM and we observe a population of projects facing the same $z$ trajectory, the population-level cancellation hazard at time $t$ is:
\begin{equation}
h^{\tau}(t) = \Pr(z_t < \bar{z}^{\tau} \mid \text{not yet cancelled, } z_{t-1}),
\end{equation}
which depends on the distance $z_t - \bar{z}^{\tau}$.

For an empirically tractable approximation, we adopt the discrete-time logit hazard specification of Chapters 6 and 7:
\begin{equation}
\mathrm{logit}\, h^{\tau}(t) = \alpha^{\tau} + \beta_z^{\tau} z_t,
\end{equation}
with $\beta_z^{\tau} < 0$ (lower $z$ raises hazard). The carbon-conditional differential is:
\begin{equation}
\mathrm{logit}\, h^B - \mathrm{logit}\, h^G = (\alpha^B - \alpha^G) + (\beta_z^B - \beta_z^G) z_t.
\end{equation}
Matching to the empirical specification of Chapter 6 with $\eta_{it} \supset \bblue \cdot \mathrm{Blue}_i + \bint \cdot \mathrm{Blue}_i \cdot z_t$, we identify:
\begin{align}
\bblue &= \alpha^B - \alpha^G, \\
\bint &= \beta_z^B - \beta_z^G.
\end{align}

% =========================================================================
\section{Comparative Statics and Empirical Predictions}
\label{sec:ch3-comparative}

\subsection{Sign of $\bint$}

The model's central prediction concerns the sign of $\bint = \beta_z^B - \beta_z^G$.

\begin{proposition}
\label{prop:ch3-sign}
If Blue projects are more sensitive to the EUA price than Green projects ($b^B > b^G > 0$) and have higher unrecoverable capital ($\hat{K}^B > \hat{K}^G$), then the carbon-conditional interaction coefficient satisfies $\bint < 0$.
\end{proposition}

\begin{proof}[Sketch of proof]
The hazard rate is a decreasing function of $z$ for both technologies; both $\beta_z^B$ and $\beta_z^G$ are negative. The slope $|\beta_z^{\tau}|$ scales with $b^{\tau}$ (Blue's profit is more sensitive to $z$, so its hazard is more sensitive to $z$). Therefore $|\beta_z^B| > |\beta_z^G|$, which combined with both being negative implies $\beta_z^B < \beta_z^G < 0$ and hence $\bint = \beta_z^B - \beta_z^G < 0$. \hfill $\square$
\end{proof}

The empirical estimate of $\bint = -1.5$ in our cleanest subsample (Chapter 7) is consistent with this prediction.

\subsection{Magnitude of $\bint$}

The magnitude of $\bint$ depends on the relative profit-sensitivity to EUA prices. A naïve calibration approach: if Blue's hydrogen revenue is approximately $30\%$ more sensitive to EUA than Green's (because Blue benefits from grey-hydrogen displacement), and the typical project margin is $\$ 1/$ kg, then a one-standard-deviation EUA shock translates to a $\sim 30\%$ relative profit advantage for Blue. Translated to a log-hazard ratio via the relationship $\beta_z^{\tau} \sim -b^{\tau} / \sigma$, this implies $\bint$ on the order of $-1$ to $-2$, consistent with our empirical $\bint = -1.5$.

\subsection{Time variation of $\bint$}

Chapter 7 documents that $\bint$ in the North American subsample has gradually strengthened from $\approx -0.7$ in 2014 to $\approx -1.9$ in 2024. The model can rationalise this in two ways:
\begin{enumerate}
    \item \textbf{Strengthening of the displacement mechanism}: as the EU ETS matured (post-2018 Market Stability Reserve, post-2022 CBAM announcement), the marginal grey-hydrogen producer became increasingly likely to exit at high EUA. This raises $b^B$ relative to $b^G$, strengthening the carbon-conditional differential.
    \item \textbf{Capital structure shifts}: as PEM electrolyser CAPEX fell from approximately \$ 1200/kg/day in 2014 to below \$ 600/kg/day in 2024 (\cite{iea2024hydrogen}), the gap $\hat{K}^B - \hat{K}^G$ widened. This increases Blue's relative cancellation sensitivity.
\end{enumerate}
Both mechanisms predict $|\bint(t)|$ increasing over time, consistent with the GAS posterior trajectory in Chapter 7.

\subsection{Robustness to alternative parameterisations}

If the assumption $b^B > b^G$ is reversed (Green benefits more from high EUA due to absence of residual carbon costs), the model predicts $\bint > 0$. Our empirical finding of $\bint < 0$ thus serves as a model selection device, supporting the competitive-displacement interpretation over the cost-pass-through interpretation. We return to this in Section \ref{sec:ch3-extensions}.

% =========================================================================
\section{Extensions and Limitations}
\label{sec:ch3-extensions}

\subsection{Multiple state variables}

The single-state model omits factors that are likely to interact with EUA in determining cancellation propensity. Gas prices (TTF) affect Blue's operating costs; electricity prices affect Green's; capital availability and financing conditions affect both. A multi-state extension would model these jointly, with the carbon-conditional effect identified by EUA-specific variation conditional on the other states. The cost is greater empirical demands; the benefit is sharper identification of the carbon channel.

\subsection{Heterogeneous projects}

The model assumes a representative Blue and Green project. In practice, projects differ in scale, sponsor type, regulatory environment, and contract structure. Our empirical heterogeneous treatment effects analysis (Chapter 7, Section 7.5) provides partial evidence that the carbon-conditional effect concentrates among certain project types, although the small sample size limits precise inference. Extensions could allow $b^{\tau}$ and $\hat{K}^{\tau}$ to vary by project type.

\subsection{Policy uncertainty}

EUA price is itself partly a function of policy uncertainty (ETS cap revisions, Market Stability Reserve activations, CBAM implementation timeline). A more sophisticated framework would model regime-switching dynamics for $\mu$ and $\sigma$ \cite{creal2013generalized}, allowing the cancellation threshold to vary across policy regimes.

\subsection{Strategic behaviour}

Project sponsors may strategically time their cancellations to influence policy formation (e.g., cancellations that signal infeasibility to motivate stronger carbon pricing). Our model assumes price-taking behaviour. An extension with sponsor-policy interaction would be of independent interest.

% =========================================================================
\section{Policy mechanism taxonomy: which friction does each instrument reduce?}
\label{sec:ch3-mechanism-taxonomy}

The real-options framework of Sections~\ref{sec:ch3-setup}--\ref{sec:ch3-comparative} characterises the optimal cancellation policy under a single source of uncertainty (the EUA price). Realistic policy environments contain multiple economic frictions simultaneously, and the central empirical observation that policy effects are heterogeneous across instruments (the carrot-policy findings of Chapter~\ref{ch:results_did}) is best understood as different instruments reducing different subsets of these frictions. This section develops an explicit taxonomy that makes the friction-to-instrument mapping operational and links it to the empirical estimates.

The framework follows the mechanism-design tradition of \cite{laffont1993theory} and the financial-frictions literature of \cite{holmstrom1997financial,holmstrom1998private}, supplemented with the directed-technical-change perspective of \cite{acemoglu2012environment}. The taxonomy is not exhaustive of the broader environmental-economics literature, but it is sufficiently complete to classify each policy instrument analysed in this thesis.

\subsection{Seven economic frictions in clean-hydrogen investment}
\label{sec:ch3-frictions}

The decision to commit capital to a hydrogen project faces seven distinct frictions, each of which can in principle be reduced by an appropriately designed policy instrument:

\paragraph{F1 --- Revenue uncertainty.} The mark-to-market revenue of the hydrogen output is unknown at the moment of FID. The buyer side of the market is thin, there are no liquid spot or futures prices for low-carbon hydrogen as of the sample window, and revenue depends on bilateral negotiation with industrial offtakers. This is the friction that production tax credits (US Section 45V, US Section 45Q) address by guaranteeing a per-unit payment indexed to output, conditional on eligibility verification.

\paragraph{F2 --- Carbon-payoff uncertainty.} The marginal economic value of clean hydrogen relative to its fossil counterpart depends on the difference between the carbon cost of fossil hydrogen and the carbon cost of clean hydrogen. In the EU ETS, this differential is driven by the EUA price and is the state variable of the Chapter~\ref{ch:theory} real-options model. Carbon-payoff uncertainty is the friction addressed by carbon-price-floor proposals, by long-term EUA auction reserve guarantees, and indirectly by border-adjustment mechanisms such as CBAM.

\paragraph{F3 --- Coordination failure.} Hydrogen requires complementary investments in production, transport, storage, and end-use facilities. A single sponsor cannot internalise the network externality that arises when multiple actors must invest simultaneously for any of them to recover the investment. Cluster-tender mechanisms (UK Track-1, UK HAR-2, Dutch HEAVENN), state-coordinated procurement (China 14th FYP), and IPCEI-style multilateral funding are the canonical instruments addressing coordination failure.

\paragraph{F4 --- Financing constraint.} \cite{holmstrom1997financial} establish that under information frictions, the price of external finance exceeds the price of internal finance, generating capital rationing for new technologies without established cashflow track records. The first-best instrument is a direct capex grant that substitutes external for internal finance one-for-one. The EU Innovation Fund, US DOE H2Hubs, and Japan's NEDO subsidies are pure F4-addressing instruments.

\paragraph{F5 --- Implementation and execution risk.} Even with secured revenue, complete coordination, and adequate financing, projects face technological-learning risk, regulatory-permit risk, and construction-cost overrun risk. The 45V three-pillars rules (additionality, hourly time matching, geographic deliverability) constitute a particular kind of implementation friction: the policy itself adds compliance complexity that materially raises F5 for eligible projects. The triple-difference result of Section~\ref{sec:results_did_45v} ($\delta_{DDD} = +0.285$) quantifies this perverse F5-amplifying effect of the December 2023 NPRM.

\paragraph{F6 --- Selection and attrition.} Subsidy auctions and tender-based allocation generate adverse selection: when the policy criterion does not perfectly align with the social-welfare criterion, winners may be projects whose private cost structure is favourable but whose social value is moderate, while higher-value projects are screened out (the \cite{myerson1981optimal} optimal-auction logic in reverse). The UK Track-1 selection-funnel evidence (Chapter~\ref{ch:results_did}, Appendix~\ref{app:case_uk_track}) is a clean illustration: the positive headline coefficient on Track-1 is consistent with the policy selecting projects that would have proceeded irrespective of the Track-1 award, rather than with Track-1 causing additional realisations.

\paragraph{F7 --- Counterparty risk and incomplete contracting.} \cite{grossman1986costs} formalise the inefficiencies that arise when contracts cannot specify all relevant contingencies. In hydrogen specifically, a long-term offtake contract may not be enforceable at the bilateral level if the offtaker's plant becomes economically obsolete or if regulatory changes alter the offtaker's compliance needs. Pre-FID offtake-commitment reduces F7 by establishing a credible counterparty arrangement before the irreversible production investment is made. The 11--13 percentage-point survival premium documented in Chapter~\ref{ch:results_offtake} is the empirical magnitude of F7-reduction by voluntary contracting.

\subsection{Friction-to-instrument mapping}
\label{sec:ch3-friction-mapping}

Table~\ref{tab:ch3-friction-instrument} maps each policy instrument analysed in this thesis to its primary and secondary frictions reduced. The matrix is rectangular by design: most real-world instruments address more than one friction, but always with a clear primary target. Reading the table down a column reveals which frictions are heavily addressed (financing, revenue) versus under-addressed (coordination, counterparty) in current hydrogen policy. Reading across a row reveals which combinations of frictions are jointly tackled by each instrument.

\begin{table}[!htbp]
\centering
\caption{Policy instruments and the economic frictions they reduce}
\label{tab:ch3-friction-instrument}
\small
\begin{tabular}{lccccccc}
\toprule
Instrument & F1 Rev & F2 Carb & F3 Coord & F4 Fin & F5 Impl & F6 Sel & F7 Counter \\
\midrule
US Section 45Q & $\bullet$ & $\bullet$ & --- & $\circ$ & --- & --- & --- \\
US Section 45V (PTC) & $\bullet$ & --- & --- & $\circ$ & --- & --- & --- \\
US 45V three-pillars NPRM & $\circ$ & --- & --- & --- & $\Box$ & --- & --- \\
EU Innovation Fund (capex grants) & --- & --- & --- & $\bullet$ & --- & $\Box$ & --- \\
EU Innovation Fund Auction (Hydrogen Bank) & $\bullet$ & --- & --- & $\bullet$ & --- & $\Box$ & --- \\
EU CBAM (transitional 2023--2025) & --- & $\bullet$ & --- & --- & --- & --- & --- \\
EU CBAM (definitive 2026+) & --- & $\bullet$ & --- & --- & --- & --- & --- \\
UK Track-1 / HAR-1 & $\bullet$ & --- & $\bullet$ & $\bullet$ & --- & $\Box$ & $\circ$ \\
UK HAR-2 (CfD-based) & $\bullet$ & --- & --- & --- & --- & $\Box$ & --- \\
China 14th Five-Year Plan & $\circ$ & --- & $\bullet$ & $\bullet$ & --- & --- & $\bullet$ \\
Pre-FID offtake commitments & $\bullet$ & --- & $\circ$ & --- & --- & --- & $\bullet$ \\
\bottomrule
\multicolumn{8}{l}{\footnotesize $\bullet$ Primary friction reduced; $\circ$ Secondary friction reduced; $\Box$ Friction \emph{amplified} by instrument design.} \\
\multicolumn{8}{l}{\footnotesize F1: Revenue uncertainty; F2: Carbon-payoff uncertainty; F3: Coordination failure; F4: Financing constraint;} \\
\multicolumn{8}{l}{\footnotesize F5: Implementation/execution risk; F6: Selection/attrition; F7: Counterparty/incomplete contracting.} \\
\end{tabular}
\end{table}

\subsection{Linking friction reduction to identified treatment effects}
\label{sec:ch3-friction-effects}

The taxonomy delivers two testable predictions that are confirmed by the empirical analysis of Chapters~\ref{ch:results_did} and~\ref{ch:results_offtake}.

The first prediction is that \emph{instruments addressing multiple binding frictions should generate larger and more robust treatment effects than instruments addressing a single friction}. The China FYP (F3+F4+F7) is the only carrot-policy with breakdown $M^* = 1.5$ in the honest-DiD sensitivity analysis (Section~\ref{sec:honest_did}); it is also the instrument reducing the largest set of frictions in Table~\ref{tab:ch3-friction-instrument}. The EU Innovation Fund (F4 only) yields the cleanest informative null (six convergent methods); it addresses the one friction least binding in the contemporary hydrogen environment, where capital is available but counterparty commitments and coordination are missing. The 45Q effect (F1+F2) is point-identified at L3.a but bounded at L3.b ($M^* = 0.2$); it addresses an intermediate set of frictions.

The second prediction is that \emph{instruments that amplify rather than reduce frictions should produce treatment effects with the opposite sign of protective effects}. The US 45V three-pillars NPRM converts the F1-reducing 45V tax credit into an F5-amplifying compliance regime via the additionality, hourly matching, and deliverability requirements. The DDD estimate of $+0.285$ in Section~\ref{sec:results_did_45v} is the empirical magnitude of this friction-amplifying effect: a 28.5 percentage point increase in US-Green hydrogen cancellation rate relative to the within-jurisdiction blue placebo.

\subsection{The offtake mechanism and the $\sigma$-channel}
\label{sec:ch3-offtake-sigma}

The Chapter~\ref{ch:results_offtake} offtake-commitment effect occupies a distinct position in the taxonomy. It is not a government policy but a contractual instrument; nonetheless it reduces F7 (counterparty risk) and partially F1 (revenue uncertainty). The cross-sectoral heterogeneity test --- the offtake-effect is concentrated in high-revenue-volatility sectors (power and heat, transport) and small or null in low-volatility sectors (chemical, industry) --- provides a direct empirical confirmation of the F1+F7 channel: offtake-commitment is most valuable where revenue volatility is highest, which is what the $\sigma$-channel of Chapter~\ref{ch:theory} predicts.

The substantive interpretation is that the offtake-commitment mechanism is a private-contracting analogue of public-policy instruments addressing the same friction pair. Where the China FYP achieves F3+F4+F7 reduction through state procurement mandates, voluntary offtake commitments achieve F1+F7 reduction through bilateral contracting. The two are partial substitutes: the China FYP $\times$ offtake interaction in Section~\ref{sec:offtake_policy_interaction} of Chapter~\ref{ch:results_offtake} ($+0.108$, $p = 0.003$) shows that the marginal value of voluntary offtake-commitment is reduced where state procurement mandates already exist.

\subsection{Implications for instrument design}
\label{sec:ch3-design-implications}

The taxonomy supports three design implications, formalised in Chapter~\ref{ch:conclusion}.

First, \emph{capex grants without complementary offtake mechanisms address only F4}, which is not the binding friction in the contemporary hydrogen environment. The empirical EU IF null is the consequence: with capital available at low cost, the marginal grant does not change the F4-shadow price meaningfully, while leaving F1+F7 unaddressed. Reform of EU IF to require demonstrable pre-FID offtake commitments would address F7 jointly with F4 and is the policy proposal of Section~11.2.

Second, \emph{output-credit mechanisms (US 45Q-style) are F1-reducers but are vulnerable to F5-amplifying compliance designs}. The 45V three-pillars NPRM is the canonical illustration of how an F1-reducing instrument can be functionally transformed into an F5-amplifier by overly stringent verification requirements. The 45V Final Rules of January 2025 partially walked back these provisions; the resulting effect-size attenuation is a natural-experiment test of the friction-classification.

Third, \emph{coordination-failure instruments (cluster tenders, state procurement) are subject to F6 selection effects} that can produce apparent positive treatment effects which on closer inspection reflect the selection process rather than additional realisations. The UK Track-1 selection-funnel analysis is the canonical illustration. Mechanism-design treatments of subsidy auctions \cite{laffont1993theory, myerson1981optimal} suggest that scoring rules favouring marginal-additionality projects rather than highest-bid projects would partially address F6.

\section{Connection to Empirical Findings}
\label{sec:ch3-empirical-connection}

The model yields three testable predictions, each addressed empirically in Chapters 6 and 7:

\begin{enumerate}
    \item \textbf{$\bint < 0$.} Empirical estimate: $\bint = -1.17$ (NA Static), $\bint = -1.88$ (NA GAS long-run mean) -- both negative and significant. \textbf{Consistent.}
    \item \textbf{$|\bint|$ increases when the capital intensity gap widens.} Empirical evidence: NA-only GAS shows $|\bint(t)|$ increasing from $\sim 0.7$ in 2014 to $\sim 1.9$ in 2024, period over which the Green-CAPEX-to-Blue-CAPEX ratio has fallen sharply. \textbf{Consistent.}
    \item \textbf{Heterogeneity by project commitment.} Empirical evidence: positive sub-coefficient (or null effect) in the EU subsample where all PEM cancellations come from ``Unknown'' sponsors (likely pilot/research projects with low commitment), versus clearly negative coefficient in the NA subsample (named major sponsors). \textbf{Consistent} with the interpretation that the mechanism is sharper among genuinely committed projects.
\end{enumerate}

This three-way consistency provides reasonable support for the structural interpretation of the empirical findings.

% =========================================================================
\section{Conclusion}
\label{sec:ch3-conclusion}

This chapter has developed a real-options framework for the cancellation decision in hydrogen technology investments under EUA-price uncertainty. The model identifies two structural drivers of the carbon-conditional cancellation differential: differential profit-sensitivity to EUA (with Blue benefiting more from high EUA due to grey-hydrogen displacement) and differential capital intensity (with Blue's higher sunk costs raising its cancellation threshold). The combination yields a testable prediction of $\bint < 0$ that is consistent with the empirical estimates of Chapters 6 and 7. The model further provides a structural rationale for the empirically observed temporal intensification of the mechanism, attributable both to maturation of the carbon pricing regime and to falling Green-technology CAPEX.

% =========================================================================
\begin{thebibliography}{99}

\bibitem[Bolton and Kacperczyk(2021)]{bolton2021carbon}
Bolton, P., \& Kacperczyk, M. (2021).
\newblock Do investors care about carbon risk?
\newblock \textit{Journal of Financial Economics}, 142(2), 517--549.

\bibitem[Creal et~al.(2013)]{creal2013generalized}
Creal, D., Koopman, S.~J., \& Lucas, A. (2013).
\newblock Generalized autoregressive score models with applications.
\newblock \textit{Journal of Applied Econometrics}, 28(5), 777--795.

\bibitem[Dixit and Pindyck(1994)]{dixit1994investment}
Dixit, A.~K., \& Pindyck, R.~S. (1994).
\newblock \textit{Investment under Uncertainty}.
\newblock Princeton University Press.

\bibitem[IEA(2024)]{iea2024hydrogen}
International Energy Agency (2024).
\newblock \textit{Global Hydrogen Review 2024}.
\newblock IEA, Paris.

\bibitem[Odenweller et~al.(2022)]{odenweller2022probabilistic}
Odenweller, A., Ueckerdt, F., Nemet, G.~F., Jensterle, M., \& Luderer, G. (2022).
\newblock Probabilistic feasibility space of scaling up green hydrogen supply.
\newblock \textit{Nature Energy}, 7(9), 854--865.

\bibitem[Pindyck(1991)]{pindyck1991irreversibility}
Pindyck, R.~S. (1991).
\newblock Irreversibility, uncertainty, and investment.
\newblock \textit{Journal of Economic Literature}, 29(3), 1110--1148.

\end{thebibliography}

